% Sources: Per Alexandersson, SymCat, 2026:
% https://www.symmetricfunctions.com/realRootedOpenProblems.htm#conjNestedMatroidZ
% https://www.symmetricfunctions.com/schubertMatroidZ.htm
% The author's equivalent question is quoted verbatim from the opening of
% https://github.com/PerAlexandersson/RealRooted/issues/2650
\documentclass{article}
\usepackage[T1]{fontenc}
\usepackage{amsmath,amssymb}
\begin{document}
\paragraph{Question (Alexandersson, 2026; verbatim).}
Are the Z-polynomials (Proudfoot--Xu--Young) of \textbf{nested matroids} real-rooted?

\paragraph{Definitions.}
On $[n]=\{1,\ldots,n\}$ the nested matroid specified by upper bounds
$q_1\le\cdots\le q_r\le n$ has family of bases
\[
 \bigl\{\{b_1,\ldots,b_r\}:1\le b_1<\cdots<b_r\le n,\quad
 b_j\le q_j\ \ (1\le j\le r)\bigr\},
\]
where the family is assumed nonempty. For a straight shape with positive
parts $\lambda_1\ge\cdots\ge\lambda_r$, use
$n=r+\lambda_1$ and $q_j=j+\lambda_{r+1-j}$.
For the skew shape $\lambda/\mu$ in the extension, impose also
$b_j\ge j+\mu_{r+1-j}$, with $\mu_r=0$.
Rows $i,i+1$ have $\lambda_{i+1}-\mu_i$ columns in common when this is
positive.

Write $\mathcal L(M)$ for the lattice of flats (sets closed under adjoining
elements without rank increase), $\rho$ for the rank function and $M/F$
for contraction. The matroid $Z$-polynomial is
\[
 Z_M(t)=\sum_{F\in\mathcal L(M)}t^{\rho(F)}P_{M/F}(t)
\]
(Proudfoot, Xu and Young, \emph{The $Z$-polynomial of a matroid},
Electronic J. Combin.\ 25(1), P1.26, doi:10.37236/7105).
In the loopless setting here, the Kazhdan--Lusztig polynomials $P$ are
equivalently characterized by $P_M=1$ for rank zero, $\deg P_M<r/2$
for positive rank $r$, and $t^r Z_M(1/t)=Z_M(t)$ for every $M$.
Real-rooted means that every complex zero is real.

\paragraph{Scope of the result.}
The accompanying proof establishes strict real-rootedness and a common
Legendre interlacer when $r\ge2$ and $\lambda_r\ge30(r-1)$, also for
skew shapes when each consecutive-row overlap is at least $30(r-1)$.
It is not a solution for all nested matroids.
\end{document}
